\documentclass[11pt]{amsart}
\usepackage[T1]{fontenc}

\usepackage{amsmath,amssymb,amsthm}
\usepackage{mathtools}
\usepackage[margin=1.15in]{geometry}
\usepackage{xcolor}
\usepackage{hyperref}
\hypersetup{
  colorlinks=true,
  linkcolor=blue!60!black,
  citecolor=blue!60!black,
  urlcolor=blue!60!black,
  pdftitle={The Taylor coefficients of the Riemann xi-function form a Polya frequency sequence of order 3},
  pdfauthor={Cyril Rodaro},
  pdfsubject={Polya frequency order of the Riemann xi-function coefficients},
  pdfkeywords={Riemann xi-function, Polya frequency sequences, total positivity, Toeplitz minors}
}

\newtheorem{theorem}{Theorem}[section]
\newtheorem{lemma}[theorem]{Lemma}
\newtheorem{proposition}[theorem]{Proposition}
\newtheorem{corollary}[theorem]{Corollary}
\theoremstyle{definition}
\newtheorem{definition}[theorem]{Definition}
\theoremstyle{remark}
\newtheorem{remark}[theorem]{Remark}

\DeclareMathOperator{\PF}{PF}
\newcommand{\Xice}{\Xi}
\newcommand{\minor}[2]{\Delta_{#1}^{#2}}

\title[The $\xi$ coefficients form a $\PF_3$ sequence]{The Taylor coefficients of the Riemann $\xi$-function\\ form a P\'olya frequency sequence of order $3$}

\author{Cyril Rodaro}
\address{Independent researcher}
\email{cyril-rodaro@outlook.fr}

\date{July 2026}

\begin{document}

\begin{abstract}
Let $\xi(s)=\tfrac12 s(s-1)\pi^{-s/2}\Gamma(s/2)\zeta(s)$ and write
$\xi(\tfrac12+z)=\sum_{n\ge 0}\gamma_n z^{2n}/n!$, so that all
$\gamma_n>0$. We prove, unconditionally, that the normalized sequence
$a_n=\gamma_n/n!$ is a P\'olya frequency sequence of order $3$: every
minor of order at most $3$ of the one-sided Toeplitz matrix
$(a_{j-i})_{i,j\ge 0}$ is nonnegative. By the theory of
Aissen--Schoenberg--Whitney--Edrei, the corresponding statement for
minors of \emph{every} order is equivalent to the Riemann Hypothesis;
order $2$ (the Tur\'an inequalities) was established by Csordas,
Norfolk and Varga in 1986, and the degree-$3$ statement on the parallel
Jensen-polynomial ladder by Dimitrov and Lucas in 2011.

The proof combines four ingredients: (i) a structure theorem reducing
all consecutive-row order-$3$ minors of a positive one-sided Toeplitz
sequence to four ratio conditions; (ii) a completion theorem, proved by
two signed Grassmann--Pl\"ucker inductions, showing that for strictly
log-concave sequences the consecutive-row case already implies full
$\PF_3$; (iii) certified interval-arithmetic verification of the ratio
conditions for $2\le n\le 126$; and (iv) an analytic proof of the ratio
conditions for all $n\ge 127$ via two-sided bounds on the cumulants of
an exponentially tilted moment family, with explicit rational
constants.

The result is sharp for the method: we exhibit explicit rational
sequences satisfying all four ratio conditions, and even full
window-$\PF_3$, for which an order-$4$ Toeplitz minor is negative.
Consequently no ratio-type hypothesis of the kind used here can reach
$\PF_4$, and further progress on the P\'olya frequency ladder for $\xi$
requires structural information specific to $\xi$. All computational
claims are backed by a reproducible artifact with fail-closed
manifests, pinned dependency hashes, and deterministic replay.
\end{abstract}

\maketitle

\section{Introduction}\label{sec:intro}

Write
\begin{equation}\label{eq:xi}
\xi(s)=\tfrac12\,s(s-1)\,\pi^{-s/2}\,\Gamma(s/2)\,\zeta(s),
\qquad
\Psi(z):=\xi(\tfrac12+z)=\sum_{n\ge 0}\frac{\gamma_n}{n!}\,z^{2n},
\end{equation}
so that $\gamma_n>0$ for all $n\ge 0$. The Riemann Hypothesis (RH)
asserts that all zeros of $\Psi$ are purely imaginary; equivalently,
setting
\begin{equation}\label{eq:F}
F(w):=\Psi(\sqrt{w})=\sum_{n\ge 0}a_n w^{n},
\qquad a_n:=\frac{\gamma_n}{n!}>0,
\end{equation}
RH asserts that the entire function $F$, which has genus $0$, has only
real nonpositive zeros.

A classical theorem of Aissen, Schoenberg and Whitney, completed by
Edrei \cite{ASW52,Edrei52}, characterizes exactly this situation in
terms of total positivity: a sequence $(a_n)_{n\ge0}$ of nonnegative
numbers is a \emph{P\'olya frequency sequence} ($\PF_\infty$), i.e.\
every finite minor of the one-sided Toeplitz matrix
\[
T(a)=(a_{j-i})_{i,j\ge 0},\qquad a_k:=0 \text{ for } k<0,
\]
is nonnegative, if and only if its generating function has the form
$C z^m e^{\gamma z}\prod_i (1+\alpha_i z)/(1-\beta_i z)$ with
$C,\gamma,\alpha_i,\beta_i\ge 0$. For the genus-$0$ function $F$ of
\eqref{eq:F} this reduces to: \emph{$(a_n)$ is $\PF_\infty$ if and only
if RH holds}. Total positivity of finite order therefore provides a
ladder of unconditional statements converging to RH: for $r\ge 1$, the
sequence is $\PF_r$ when every minor of $T(a)$ of order at most $r$ is
nonnegative.
Classical sources and modern treatments of this circle of ideas include
\cite{Fekete13,Hutchinson23,Karlin68,Pinkus10,PolyaCollected}.

The history of this ladder, and of its close relative built from Jensen
polynomials, is short. Order $2$ --- log-concavity of $(a_n)$, which
amounts to $\gamma_n^{2}\ge\frac{n}{n+1}\gamma_{n-1}\gamma_{n+1}$ and
is therefore implied by (and slightly weaker than) the Tur\'an
inequalities $\gamma_n^{2}\ge\gamma_{n-1}\gamma_{n+1}$ --- follows from
the work of Csordas, Norfolk and Varga \cite{CNV86}. On the
Jensen ladder, hyperbolicity of the degree-$3$ Jensen polynomials for
all shifts (the ``higher order Tur\'an inequalities'') was proved by
Dimitrov and Lucas \cite{DL11}. Griffin, Ono, Rolen and Zagier
\cite{GORZ19} (see also \cite{GORTTW,OSullivan}) later proved that for
each fixed degree $d$ the Jensen polynomials $J^{d,n}$ are hyperbolic
for all sufficiently large shifts $n$. Families of \emph{contiguous}
order-$3$ determinants attached to $\xi$ also appear in the
determinantal programme of Nuttall \cite{Nuttall11}, which discusses
positivity for particular contiguous subfamilies. To our knowledge,
however, no unconditional statement covering \emph{all} order-$3$
Toeplitz minors of $(a_n)$ --- sparse rows and columns included ---
has appeared. The present paper closes this order:

\begin{theorem}[Main Theorem]\label{thm:main}
The sequence $a_n=\gamma_n/n!$ of normalized Taylor coefficients of
$\xi$ is $\PF_3$: every minor of order at most $3$ of the one-sided
Toeplitz matrix $T(a)$ is nonnegative. Moreover the order-$1$ entries
and all nonzero order-$2$ minors are strictly positive, and the
contiguous order-$3$ minors are strictly positive.
\end{theorem}

The proof is organized around two abstract theorems and two
$\xi$-specific inputs.

\begin{itemize}
\item[\textbf{(A)}] A \emph{structure theorem}
(Theorem~\ref{thm:consec}): for any positive one-sided Toeplitz
sequence, four ratio conditions --- log-concavity, monotonicity of the
second-order ratios $q_n$, the initial bound $q_2\le\tfrac12$, and
nonnegativity of the \emph{contiguous} order-$3$ minors --- imply
nonnegativity of every order-$3$ minor with three \emph{consecutive
rows} (arbitrary columns).
\item[\textbf{(B)}] A \emph{completion theorem}
(Theorem~\ref{thm:plucker}): for strictly log-concave positive
sequences, the consecutive-row case of order $3$ already implies full
$\PF_3$. The proof runs two signed three-term Grassmann--Pl\"ucker
inductions in which every divisor is strictly positive and no
sign-unknown minor ever appears with a negative coefficient.
\item[\textbf{(C)}] \emph{Certified finite verification}: the four
ratio conditions hold for $2\le n\le 126$, verified in ball arithmetic
(Arb \cite{Johansson17}) with rigorous enclosures of $\gamma_2,\dots,
\gamma_{130}$ and directed rounding throughout.
\item[\textbf{(D)}] An \emph{analytic tail}: for all $n\ge 127$ the
ratio conditions follow from new two-sided bounds on the second and
third cumulants of an exponentially tilted moment family attached to
the classical integral representation of $\Xi$, with explicit rational
constants; the conversion to the discrete conditions is a finite list
of coefficient-positive polynomial identities.
\end{itemize}

The result is \emph{sharp for its method} in a strong sense.

\begin{theorem}[Sharpness; see \S\ref{sec:sharp}]\label{thm:sharp-intro}
There exist explicit rational positive sequences satisfying all four
ratio conditions of (A) --- indeed satisfying every order-$3$
nonnegativity requirement on a window containing the failure --- for
which an explicit order-$4$ Toeplitz minor is negative. There also
exists a positive sequence with entire generating function satisfying
the same ratio conditions for which already the degree-$3$ Jensen
polynomial $J^{3,0}$ fails to be hyperbolic.
\end{theorem}

Thus ratio-type conditions of the kind that power Theorem~\ref{thm:main}
cannot prove $\PF_4$, nor even degree-$3$ Jensen hyperbolicity, for any
abstract sequence: any ascent of the ladder beyond order $3$ must use
information specific to $\xi$ that is invisible to these invariants.
We regard this negative statement as of equal practical importance to
the positive one; it delimits precisely where the present technique
stops. Under RH, of course, $(a_n)$ is $\PF_r$ for every $r$; our
countermodels are abstract sequences, not related to $\xi$.

\subsection*{Relation to the Jensen ladder}
Theorem~\ref{thm:main} and the Dimitrov--Lucas theorem \cite{DL11} are
statements on two different ladders which intersect only at infinity.
Hyperbolicity of degree-$d$ Jensen polynomials is a discriminant-type
condition; $\PF_r$ is a determinantal one. At finite level neither
family of statements implies the other in any obvious way, and the
countermodels of \S\ref{sec:sharp} separate them explicitly in one
direction. The two ladders coincide, by \cite{ASW52,Edrei52} and by
P\'olya's classical equivalence, only in the limit
$r=\infty$/$d=\infty$, where both are RH.

\subsection*{Computer assistance and reproducibility}
Steps (C) and parts of (D) are computer-assisted in the tradition of
\cite{PlattTrudgian}: every numerical claim entering the proof is a
certified ball-arithmetic enclosure with outward rounding, produced
with Arb through the open-source \texttt{python-flint} bindings and
assembled by a
fail-closed checker that pins the exact inputs by cryptographic hash
and refuses to report success unless every dependency replays. The
complete artifact (code, certificates, manifests, replay instructions)
accompanies the paper; see \S\ref{sec:comp}. The research programme
that produced the mathematics was conducted by autonomous AI reasoning
agents under human direction and audit; we describe the workflow and
the verification discipline in \S\ref{sec:disclosure}, and emphasize
that the proof as presented is human-readable and independent of any
claim about how it was found.

\section{Notation and main results}\label{sec:notation}

Throughout, $(a_n)_{n\ge 0}$ is a sequence with $a_n>0$, extended by
$a_k=0$ for $k<0$. Define
\begin{equation}\label{eq:ratios}
R_n=\frac{a_n}{a_{n-1}}\ (n\ge 1),
\qquad
q_n=\frac{R_n}{R_{n-1}}\ (n\ge 2).
\end{equation}
Log-concavity of $(a_n)$ ($\PF_2$) is $q_n\le 1$ for all $n\ge2$;
\emph{strict} $\PF_2$ is $q_n<1$, i.e.\ $R_n$ strictly decreasing.

For increasing index triples $R=(r_0,r_1,r_2)$, $C=(c_0,c_1,c_2)$ set
\[
\minor{R}{C}=\det\bigl(a_{c_j-r_i}\bigr)_{0\le i,j\le 2}.
\]
A minor is \emph{contiguous} when both $R$ and $C$ are triples of
consecutive integers; it has \emph{consecutive rows} when $R$ is. For
$m\ge 0$ we write
\[
D_m:=\minor{(0,1,2)}{(m,m+1,m+2)}
\]
for the \emph{translated contiguous determinant at start $m$}; by the
Toeplitz structure every contiguous minor is some $D_m$.

\begin{definition}\label{def:ratio-conditions}
We say $(a_n)$ satisfies the \emph{ratio conditions} when:
\begin{enumerate}
\item[(R1)] $q_n\le 1$ for all $n\ge 2$ \hfill($\PF_2$);
\item[(R2)] $q_{n+1}\ge q_n$ for all $n\ge 2$;
\item[(R3)] $q_2\le \tfrac12$;
\item[(R4)] every contiguous order-$3$ minor is nonnegative.
\end{enumerate}
\end{definition}

\begin{theorem}[Consecutive rows from ratio conditions]\label{thm:consec}
If $(a_n)$ satisfies the ratio conditions, then every order-$3$ minor
of $T(a)$ with consecutive rows is nonnegative.
\end{theorem}

\begin{theorem}[Pl\"ucker completion]\label{thm:plucker}
Let $(a_n)$ be positive and strictly $\PF_2$, and suppose every
order-$3$ minor of $T(a)$ with consecutive rows is nonnegative. Then
$(a_n)$ is $\PF_3$.
\end{theorem}

\begin{theorem}[Analytic and certified inputs for $\xi$]\label{thm:inputs}
For $a_n=\gamma_n/n!$:
\begin{enumerate}
\item[(i)] $a_n>0$ for all $n\ge0$;
\item[(ii)] $q_n<1$ for all $n\ge 2$ (strict $\PF_2$);
\item[(iii)] $q_2<\tfrac12$;
\item[(iv)] $q_{n+1}\ge q_n$ for all $n\ge2$;
\item[(v)] every contiguous order-$3$ minor is strictly positive.
\end{enumerate}
\end{theorem}

Theorem~\ref{thm:main} follows at once:
Theorem~\ref{thm:inputs} verifies the ratio conditions for $\xi$,
Theorem~\ref{thm:consec} yields all consecutive-row order-$3$ minors,
and Theorem~\ref{thm:plucker} completes to $\PF_3$.

\section{Consecutive rows from the ratio conditions}\label{sec:consec}

Fix rows $(m,m+1,m+2)$; by the Toeplitz structure it suffices to treat
$m=0$, replacing columns $(c_0,c_1,c_2)$ by their translates. Columns
containing an index $<0$ produce structurally triangular or zero
minors: if $c_0<0$ the first column has at most one nonzero entry and
the minor reduces, with nonnegative cofactor by $\PF_2$, to an
order-$2$ case. So assume $0\le c_0<c_1<c_2$ and write the columns as
$(c_0,c_1,c_2)$.

The proof splits into three families according to $c_0$.

\subsection{Columns \texorpdfstring{$(0,j,k)$}{(0,j,k)}}
An exact factorization gives
\begin{equation}\label{eq:fam0}
\minor{(0,1,2)}{(0,j,k)}
= a_0\,a_{j-2}\,a_{k-2}\,\bigl(R_{j-1}-R_{k-1}\bigr),
\qquad 2\le j<k,
\end{equation}
which is nonnegative by (R1) since $j-1<k-1$. (For $j\le 2$ the minor
is triangular.)

\subsection{Columns \texorpdfstring{$(1,j,k)$}{(1,j,k)}}
With $G_c:=R_{c-1}\,(R_1-R_c)$ one has the exact factorization
\begin{equation}\label{eq:fam1}
\minor{(0,1,2)}{(1,j,k)}
= a_0\,a_{j-2}\,a_{k-2}\,\bigl(G_j-G_k\bigr),
\qquad 2\le j<k .
\end{equation}
Monotonicity of $c\mapsto G_c$ on $c\ge 2$ is exactly where (R2) and
(R3) enter: writing $G_{c+1}-G_c$ in terms of the $q$'s, the
anchored inequality $q_2\le\tfrac12$ together with $q$ nondecreasing
forces $G$ nonincreasing, hence \eqref{eq:fam1} is nonnegative.

\subsection{Columns \texorpdfstring{$(i,j,k)$, $i\ge 2$}{(i,j,k), i >= 2}}
Here the minor factors as a positive product times the orientation
determinant of the three planar points
\[
P_c=\bigl(R_{c-1},\,R_{c-1}R_c\bigr),\qquad c\in\{i,j,k\},
\]
taken in decreasing abscissa order (by strict monotonicity of
$R$). Nonnegativity of all such orientations is a discrete convexity
statement for the polygonal curve $c\mapsto P_c$; the contiguous case
$(i,i+1,i+2)$ is precisely condition (R4), and the general case follows
from the contiguous one because convexity of consecutive triples of a
planar polygonal chain with monotone abscissae implies convexity of all
triples.

A complete symbolic expansion of the three families, including all
boundary and triangular cases, is replayed in the artifact
(\S\ref{sec:comp}); the identities \eqref{eq:fam0}--\eqref{eq:fam1} are
elementary to verify by direct expansion. \qed

\begin{remark}
Conditions (R1)--(R4) are close to minimal for
Theorem~\ref{thm:consec}: dropping (R3), or replacing (R2) by any
bounded number of its instances, admits counterexamples already at
order $3$; see the artifact for exact witnesses. The sharpness examples
of \S\ref{sec:sharp} show the conclusion cannot be strengthened to
$\PF_4$.
\end{remark}

\section{The \texorpdfstring{Pl\"ucker}{Plucker} completion}\label{sec:plucker}

We prove Theorem~\ref{thm:plucker}. All identities keep the column
triple $C=(c_0,c_1,c_2)$ fixed and vary rows; $\minor{(x,y,z)}{}$
abbreviates $\minor{(x,y,z)}{C}$.

\subsection{Strict order-2 positivity for arbitrary gaps}
For rows $i<j$ and columns $p<q$ the order-$2$ minor is
$a_{p-i}a_{q-j}-a_{q-i}a_{p-j}$. If $i\le p<j\le q$ it is triangular
with strictly positive diagonal. In the interior case $p\ge j$, put
$x=p-j$, $\alpha=j-i$, $\beta=q-p$; then
\begin{equation}\label{eq:pf2gap}
a_{x+\alpha}a_{x+\beta}-a_{x+\alpha+\beta}a_{x}
= a_x a_{x+\beta}\Bigl(\ \prod_{t=1}^{\alpha}R_{x+t}
\;-\;\prod_{t=1}^{\alpha}R_{x+\beta+t}\Bigr)>0,
\end{equation}
strictly, term by term, because $R$ is strictly decreasing. Thus every
order-$2$ minor is nonnegative, and strictly positive unless it is
structurally zero.

\subsection{The auxiliary row \texorpdfstring{$d=c_1+1$}{d = c1 + 1}}
For any rows $i<j<d$ with $d:=c_1+1$, the row of $T(a)$ indexed by $d$
restricted to the columns $C$ is $(0,0,a_{c_2-c_1-1})$, whence the
exact factorization
\begin{equation}\label{eq:drow}
\minor{(i,j,d)}{}
= a_{c_2-c_1-1}\,
\det\begin{pmatrix} a_{c_0-i}&a_{c_1-i}\\ a_{c_0-j}&a_{c_1-j}
\end{pmatrix},
\end{equation}
whose scalar prefix is strictly positive (equal to $a_0$ at the edge
$c_2=d$) and whose determinant is an order-$2$ minor, nonnegative by
\eqref{eq:pf2gap}.

\subsection{Boundary reductions}
Three one-sided cases are removed before any division: if $c_2<r_2$ the
last row vanishes; if $c_1<r_2\le c_2$, expansion along the last row
factors the minor as a positive entry times an order-$2$ minor; if
$c_0<r_1\le c_1$, expansion along the first column does the same. The
remaining \emph{inductive region} is
\begin{equation}\label{eq:region}
c_1\ge r_2\quad\text{and}\quad c_0\ge r_1 ,
\end{equation}
in which $d=c_1+1$ exceeds every row in play and all divisors below are
strictly positive interior order-$2$ cases of \eqref{eq:pf2gap} (or the
triangular edge case).

\subsection{Adjacent-pair rows}
Let $R=(r,r+1,s)$ with $s>r+2$. The three-term Grassmann--Pl\"ucker
identity on the five rows $r,r+1,r+2,s,d$ reads
\begin{equation}\label{eq:pl5}
\minor{(r,r+1,s)}{}\ \minor{(r+1,r+2,d)}{}
= \minor{(r,r+1,r+2)}{}\ \minor{(r+1,s,d)}{}
+ \minor{(r,r+1,d)}{}\ \minor{(r+1,r+2,s)}{} .
\end{equation}
On the right, the first factor has consecutive rows (nonnegative by
hypothesis), the two $d$-minors are nonnegative by \eqref{eq:drow}, and
the last factor has the same adjacent pair with row gap smaller by one;
the divisor $\minor{(r+1,r+2,d)}{}$ is strictly positive in the region
\eqref{eq:region}. Induction from the consecutive base $s=r+2$ shows
every minor with an adjacent \emph{leading} pair is nonnegative. The
mirrored identity
\begin{equation}\label{eq:pl6}
\minor{(r,s-1,s)}{}\ \minor{(s-2,s-1,d)}{}
= \minor{(r,s-2,s-1)}{}\ \minor{(s-1,s,d)}{}
+ \minor{(r,s-1,d)}{}\ \minor{(s-2,s-1,s)}{}
\end{equation}
handles an adjacent \emph{trailing} pair by induction on the first row
gap. Hence every order-$3$ minor containing at least one adjacent pair
of rows is nonnegative. In both identities no sign-unknown minor occurs
with a negative coefficient and no possibly-zero minor is divided out.

\subsection{Sparse rows}
Let $R=(r,u,s)$ with $r+1<u<s$. Pl\"ucker on the five rows
$r,r+1,u,s,d$ with common row $u$ gives, with the exact orientation,
\begin{equation}\label{eq:pl7}
\minor{(r,u,s)}{}\ \minor{(r+1,u,d)}{}
= \minor{(r,r+1,u)}{}\ \minor{(u,s,d)}{}
+ \minor{(r,u,d)}{}\ \minor{(r+1,u,s)}{} .
\end{equation}
The divisor $\minor{(r+1,u,d)}{}$ is strictly positive in the region
\eqref{eq:region}; the two $d$-minors are nonnegative by
\eqref{eq:drow}; $\minor{(r,r+1,u)}{}$ is covered by the adjacent-pair
case; and $\minor{(r+1,u,s)}{}$ has first gap $u-(r+1)<u-r$. Induction
on $u-r$, with the adjacent-pair case as base, proves every order-$3$
minor nonnegative. Together with positivity of the entries and
\eqref{eq:pf2gap}, $(a_n)$ is $\PF_3$. \qed

\begin{remark}[Duality check]
For $L\ge\max(r_2,c_2)$ the substitution $R^*=(L-c_2,L-c_1,L-c_0)$,
$C^*=(L-r_2,L-r_1,L-r_0)$ satisfies
$\minor{R}{C}=\minor{R^*}{C^*}$ (reflection--transposition of the
Toeplitz structure). It maps adjacent column pairs to adjacent row
pairs and provides an independent cross-check of the intermediate
reduction, although it is not needed once \eqref{eq:pl7} is in place.
\end{remark}

\begin{remark}
The identities \eqref{eq:pl5}--\eqref{eq:pl7} were verified exactly on
$500$ random integer matrices each, in addition to the symbolic
expansion replayed in the artifact.
\end{remark}

\section{The inputs for \texorpdfstring{$\xi$}{xi}: certified range and analytic tail}
\label{sec:inputs}

We now prove Theorem~\ref{thm:inputs}. Recall the classical
representation (see e.g.\ \cite{PolyaCollected,CNV86})
\begin{equation}\label{eq:moment}
\xi(\tfrac12+z)=\int_0^\infty \Phi(u)\cosh(uz)\,du,
\qquad
\Phi(u)=4\sum_{m\ge1}\bigl(2\pi^2m^4e^{9u/2}-3\pi m^2e^{5u/2}\bigr)
e^{-\pi m^2 e^{2u}} ,
\end{equation}
with $\Phi>0$ (each summand is positive since
$2\pi m^2e^{2u}/3>1$). Hence
\begin{equation}\label{eq:an-moment}
a_n=\frac{\gamma_n}{n!}=\frac{M_{2n}}{(2n)!},
\qquad
M_t:=\int_0^\infty u^{t}\,\Phi(u)\,du ,
\end{equation}
so $(M_t)_{t>0}$ is the moment family of a fixed positive density and
in particular $a_n>0$, proving (i).

\subsection{Certified range \texorpdfstring{$2\le n\le 126$}{2 <= n <= 126}}
Rigorous enclosures of $\gamma_2,\dots,\gamma_{130}$ are computed in
ball arithmetic with directed outward rounding, using a
discretization of \eqref{eq:moment} whose aliasing error is bounded
rigorously (an earlier enclosure set lacking the alias bound was
detected by internal audit and repaired; the certified set used here
carries the explicit alias term). From these balls, the quantities
$q_n$, $q_{n+1}-q_n$, and every contiguous order-$3$ determinant are
evaluated with certified comparisons, establishing (ii)--(v) in the
range $2\le n\le 126$; in particular
\[
q_2=0.4651838096951545\ldots<\tfrac12 .
\]
All comparisons are fail-closed: an undecidable ball comparison aborts
the run rather than reporting success.

\subsection{The tilted cumulant bounds}
For the tail we study the cumulant generating function in the
\emph{shift} variable,
\[
\kappa(t):=\log\int_0^\infty u^{2t}\,\Phi(u)\,du ,
\qquad\text{so that}\qquad
\log a_n=\kappa(n)-\log(2n)! .
\]
Writing $\int_0^\infty u^{2t}\Phi(u)\,du=\int_{\mathbb R}e^{2ts}V(s)\,ds$
with $V(s)=e^{s}\Phi(e^{s})>0$ (substitution $u=e^{s}$), the function
$\kappa$ is, up to the reparametrization $t\mapsto2t$, the cumulant
generating function of an exponentially tilted family; in particular
$\kappa''(t)>0$ (log-convexity) and $\kappa'''(t)$ is a positive
multiple of the third cumulant of the tilt. All constants below are
stated in this normalization, which is the one used by the artifact.

\begin{proposition}[Two-sided cumulant bounds]\label{prop:cumulant}
There is an explicit threshold
$T_{71}=\bigl[71\pi e^{71}-\tfrac94\cdot 71-1\bigr]/2
\approx 7.63\cdot 10^{32}$ such that:
\begin{enumerate}
\item[(i)] $0<\kappa''(t)<\dfrac{4}{5t}$ for all $t\ge 125$;
\item[(ii)] $\kappa'''(t)<\dfrac{11}{20\,t^{2}}$ for all $t\ge 125$;
\item[(iii)] $t^{2}\kappa'''(t)\ge -\dfrac34$ for $125\le t\le T_{71}$,
and $t^{2}\kappa'''(t)\ge -\dfrac{141}{142}$ for all $t\ge T_{71}$.
In particular $\kappa'''(t)\ge -\frac{141}{142}\,t^{-2}$, with constant
strictly below $1$, for every $t\ge 125$.
\end{enumerate}
\end{proposition}

We stress the logical roles: the \emph{lower} bound (iii) on the third
cumulant is what drives the monotonicity $q_{n+1}\ge q_n$ (it prevents
the third difference of $\log a_n$ from going negative), while the
upper bounds (i)--(ii) drive $q_n<1$ and the positivity of the
contiguous determinants. A one-sided upper bound on $\kappa'''$ alone
would be vacuous for monotonicity.

The proof of (i)--(ii) and of (iii) on the compact and banded ranges
assembles $552$ rigorous Arb boxes with directed endpoints, together
with analytic band certificates and exact rational endpoint envelopes;
the worst certified box values are
\[
\text{(second moment)}\quad 1.2372479\ldots<\tfrac54,
\qquad
t^2\kappa'''(t)\le 0.1248013790\ldots<\tfrac18
\]
on the compact part, and from the exact endpoint regime onward the
envelopes, e.g.
\[
t\,\kappa''(t)\ \le\
\tfrac{46606678326443294424723416932878}
{492590289071124220328797309123325},
\]
decrease monotonically. The bound (iii) beyond $T_{71}$ is an
\emph{effective transfer estimate} from the dominant saddle family of
the kernel to the full kernel: with $r=r(t)$ the dominant saddle
parameter defined by $\pi r e^{r}=2t+\tfrac94 r+1$, one proves
\[
\kappa'''(t)\ \ge\ -\Bigl[\frac{70}{r}+\frac{45450578214}{t}\Bigr]
\frac{1}{t^{2}}\qquad (r\ge\tfrac52),
\]
whose bracket is $\le\frac{141}{142}$ once $r\ge 71$, i.e.\ $t\ge
T_{71}$; the intermediate range $125\le t\le T_{71}$ is covered by the
certified boxes and the analytic band certificates (with the stronger
constant $\tfrac34$). Full details, including the exact rational
constants and their monotonicity certificates, are in the artifact.

\subsection{From cumulant bounds to the discrete conditions}
Since $\log a_n=\kappa(n)-\log(2n)!$, second and third finite
differences of $\log a_n$ are, by the exact integral representations
\[
\Delta^2\kappa(m)
=\iint_{[0,1]^2}\kappa''\bigl(m+x+y\bigr)\,dx\,dy,
\qquad
\Delta^3\kappa(m)
=\iiint_{[0,1]^3}\kappa'''\bigl(m+x{+}y{+}z\bigr)\,dx\,dy\,dz,
\]
controlled two-sidedly by Proposition~\ref{prop:cumulant} (all
arguments lie in $[n-2,n+1]\subset[125,\infty)$ for $n\ge 127$), while
the factorial parts are explicit rationals. This yields, for
$n\ge127$, explicit \emph{rational} upper envelopes $q_n\le X_n<1$ and
$q_{n+1}/q_n\le S_n$ (Appendix~\ref{app:envelopes}), and with
$Y_n:=S_nX_n\ (\ge q_{n+1})$ the inequalities
\[
X_n<1,\qquad Y_n<1,\qquad H(X_n,Y_n):=1-Y_n-Y_n^{2}(1-X_n)>0
\]
are proved by exhibiting coefficient-positive polynomials in
$m=n-127\ge 0$. The first gives strict $\PF_2$ in the tail. The third
enters through an exact factorization of the contiguous determinant
together with a \emph{coupled monotonicity} lemma which is what makes
the substitution of upper envelopes into $H$ legitimate
(Appendix~\ref{app:envelopes}, Lemmas~\ref{lem:A3}--\ref{lem:A5});
combined with $q$-monotonicity ($q_{n+2}\ge q_{n+1}$) this gives
$D_n>0$ for all $n\ge127$.

The monotonicity $q_{n+1}\ge q_n$ in the tail is precisely where the
\emph{lower} bound of Proposition~\ref{prop:cumulant}(iii) enters.
Indeed
\begin{align*}
\log q_{n+1}-\log q_n
  &=\Delta^{3}\kappa(n-2)-B_n,\\
B_n
  &=\log\frac{(2n+2)(2n+1)(2n-2)(2n-3)}
    {\bigl((2n)(2n-1)\bigr)^{2}}\\
  &=-\frac{2}{(n-2)^{2}}+O(n^{-3}).
\end{align*}
with $B_n$ the logarithm of an explicit rational, and the cube
representation together with (iii) bounds the cumulant term below by
$-\frac{141}{142}\,(n-2)^{-2}$ (with the stronger constant $\frac34$
below $T_{71}$), while $-B_n=2(n-2)^{-2}+O(n^{-3})$: the margin is
essentially a factor of two, and the exact rational bookkeeping shows
$\log q_{n+1}-\log q_n>0$ for every $n\ge 127$. An upper bound on
$\kappa'''$ alone would say nothing here. At the seam, the certified
enclosures extend through $\gamma_{130}$ and certify $q_n<1$ up to
$n=128$, so for the ratio bounds the certified and analytic ranges
genuinely overlap at $n\in\{127,128\}$ and agree there; for the
contiguous determinants the certified range $2\le n\le 126$ and the
analytic range $n\ge 127$ abut. Starts $m=0,1$ of the
contiguous determinants are handled explicitly ($m=0$ is triangular;
$m=1$ follows from $q_2<\tfrac12$ and $q_3>0$).

This proves (ii)--(v) for all $n$, completing
Theorem~\ref{thm:inputs}, and with it Theorem~\ref{thm:main}. \qed

\section{Sharpness}\label{sec:sharp}

The following two propositions were verified in exact rational
arithmetic; they are reproduced by a short independent script in the
artifact, and were re-verified independently by the independent AI auditor
layer.

\begin{proposition}\label{prop:cm1}
Define $(a_n)_{0\le n\le 8}$ by $a_0=1$, $R_1=1$, $R_n=R_{n-1}q_n$ and
\[
(q_2,\dots,q_8)=\Bigl(\tfrac{9}{32},\ \tfrac{19}{32},\ \tfrac34,\
\tfrac{25}{32},\ \tfrac{51}{64},\ \tfrac{15}{16},\ \tfrac{63}{64}\Bigr).
\]
Then $q$ is nondecreasing with $q_2\le\tfrac12$, and \emph{every}
order-$3$ minor of $T(a)$ with rows and columns in $\{0,\dots,8\}$ is
nonnegative; yet
\[
\det\bigl(a_{c_j-r_i}\bigr)_{\substack{r\in(0,1,2,3)\\ c\in(2,4,5,6)}}
=-\frac{809511725379305979099}{19807040628566084398385987584}<0 .
\]
Hence window-$\PF_3$ together with all ratio conditions does not imply
$\PF_4$.
\end{proposition}

\begin{proposition}\label{prop:cm2}
Define $q_2=\tfrac12$ and $q_n=\tfrac{2n-3}{2n}$ for $n\ge3$, with
$a_0=1$, $R_1=1$. The resulting positive sequence has entire
generating function, satisfies all ratio conditions (R1)--(R4), and yet
its degree-$3$ Jensen polynomial $J^{3,0}$ has discriminant
$-\tfrac{27}{16}<0$ (so is not hyperbolic) and the order-$4$ minor with
rows $(0,1,2,3)$ and columns $(1,2,3,4)$ equals $-\tfrac{5}{256}<0$.
The simpler sequence $b_n=2^{-n(n-1)/2}$ ($q_n\equiv\tfrac12$)
likewise satisfies (R1)--(R4) and has an order-$4$ minor equal to
$-\tfrac{1}{64}$.
\end{proposition}

\begin{corollary}
No consequence of the ratio conditions (R1)--(R4) --- nor of
window-verified $\PF_3$ itself --- can prove $\PF_4$, nor degree-$3$
Jensen hyperbolicity, for an abstract sequence. Any extension of
Theorem~\ref{thm:main} up the ladder must use structural information
about $\xi$ beyond these invariants.
\end{corollary}

\begin{remark}
This corollary sharpens, at the level of exact finite certificates,
the general expectation that generic positivity hypotheses cannot
climb the P\'olya frequency ladder. It is also consistent with the
observation of Farmer \cite{Farmer} that equivalences of RH tend to
concentrate, rather than dilute, the original difficulty. For kernel
--- rather than sequence --- total positivity in this landscape, see
\cite{Michalowski26}, where the de Bruijn--Newman kernel is certified
not to be $\mathrm{PF}_5$ although its low orders hold.
\end{remark}

\section{Computational methodology and reproducibility}\label{sec:comp}

All computations enter the proof only through certified interval
(ball) arithmetic with outward rounding, performed with Arb
\cite{Johansson17} through \texttt{python-flint}. The assembly layer
enforces three disciplines:

\begin{enumerate}
\item \emph{Fail-closed manifests.} Each proof stage is driven by a
checker that recomputes its claims and refuses to emit its manifest
unless every comparison is certified; undecidable comparisons abort.
\item \emph{Pinned dependencies.} A stage consumes earlier stages only
through files pinned by SHA-256 hash; the final assembly re-verifies
the full hash chain, so no stale or edited intermediate can enter
silently.
\item \emph{Deterministic replay.} Every stage exposes a single
command that replays its verification from the pinned inputs;
the top-level replay reproduces the complete proof of
Theorem~\ref{thm:inputs} and the symbolic expansions behind
Theorems~\ref{thm:consec} and~\ref{thm:plucker}.
\end{enumerate}

The artifact further contains: the exact rational countermodels of
\S\ref{sec:sharp} with an independent $30$-line verifier; the $500$
random exact tests of each Pl\"ucker identity; and the audit note
documenting the detection and repair of the aliasing gap in an earlier
enclosure set, which we include deliberately as part of the
verification history.

\smallskip
\noindent\textbf{Artifact.} The frozen supplementary archive
\texttt{xi-pf3-artifact-v1.zip} contains the proof engines,
certificates, pinned SHA-256 manifest, and a one-command fail-closed
replay. It is publicly available at
\url{https://github.com/wulf758/riemann-xi-pf3} and is frozen at commit
\href{https://github.com/wulf758/riemann-xi-pf3/commit/18aa2da35cd0f45f38041365080d76a8b32ad052}{\texttt{18aa2da35cd0f45f38041365080d76a8b32ad052}}.
The SHA-256 digest of the deposited archive is
\texttt{5D256A2580126EF4872668EBF513B431CD9861D5461786FD75F2A4D4ED9D96C0}.

\section{Discussion}\label{sec:discussion}

Theorem~\ref{thm:main} places the first unconditional rung on the
Toeplitz side of the P\'olya frequency ladder for $\xi$ beyond
log-concavity, forty years after \cite{CNV86} and fifteen after the
degree-$3$ Jensen statement of \cite{DL11}. We stress what it does
\emph{not} do: by \S\ref{sec:sharp}, the method is exactly calibrated
at order $3$, and no statement about $\PF_4$, higher Jensen degrees, or
RH follows. The countermodels quantify the obstruction: the ratio
invariants $(R_n,q_n)$, and even full window positivity at order $3$,
carry no information about order $4$.

What would be needed next is a preservation mechanism special to
$\xi$. One concrete formulation, which we leave as a problem, arises
from the exact Laguerre representation of the Jensen polynomials of
$\xi$ as positive mixtures of hyperbolic Laguerre polynomials against
the de Bruijn--Newman density: generic positive mixing does not
preserve hyperbolicity (a two-atom counterexample exists already in
degree $2$), so any proof must identify a rigidity of this specific
density (for the de Bruijn--Newman constant itself, see
\cite{RodgersTao}) --- while \cite{Michalowski26} shows that kernel-level total
positivity of the same density fails at order $5$ and therefore cannot
be that rigidity. The difficulty of RH, viewed from this ladder, is
thereby localized rather than diminished.

\section{Disclosure and acknowledgements}\label{sec:disclosure}

The mathematics in this paper was produced in an extended research
programme carried out by autonomous large-language-model agents
operating over a shared persistent memory, under human direction, with
an independent AI auditor layer; every mathematical claim was
subsequently reduced to either a human-readable proof (as presented
above), an exact rational certificate, or a certified ball-arithmetic
computation with deterministic replay. The human author directed the
programme, made all publication decisions, and takes full
responsibility for the correctness of the content. The complete
research log (some 1,300 hypothesis records, including the failed
routes and the exact countermodels that redirected the programme) is
preserved and available on request.

\appendix

\section{Explicit tail envelopes}\label{app:envelopes}

We work in the normalization of \S\ref{sec:inputs}:
$\kappa(t)=\log\int_0^\infty u^{2t}\Phi(u)\,du$ and
$\log a_n=\kappa(n)-\log(2n)!$. Throughout $n\ge 127$, so every
cumulant argument lies in $[n-2,\,n+1]\subset[125,\infty)$. The only
computer-dependent ingredients of the tail are
Proposition~\ref{prop:cumulant} itself and the single certificate of
Lemma~\ref{lem:A5}; Lemmas~\ref{lem:A1}--\ref{lem:A4} are
self-contained.

\subsection*{Envelopes}
Set
\[
f_n:=\frac{(n-1)(2n-3)}{n(2n-1)},\qquad
u_n:=\frac{4}{5(n-2)},\qquad
v_n:=\frac{11}{20(n-2)^{2}},
\]
so that $\log q_n=\Delta^{2}\kappa(n-2)+\log f_n$ and
$\log(q_{n+1}/q_n)=\Delta^{3}\kappa(n-2)+\log(f_{n+1}/f_n)$.
Proposition~\ref{prop:cumulant}(i)--(ii) give
$\Delta^{2}\kappa(n-2)\le u_n$ and $\Delta^{3}\kappa(n-2)\le v_n$, and
with $e^{w}\le 1/(1-w)$ for $0\le w<1$ we obtain the \emph{rational}
envelopes
\[
q_n\ \le\ X_n:=\frac{f_n}{1-u_n},
\qquad
\frac{q_{n+1}}{q_n}\ \le\ S_n:=\frac{f_{n+1}/f_n}{1-v_n},
\qquad
q_{n+1}\ \le\ Y_n:=S_nX_n .
\]

\begin{lemma}\label{lem:A1}
$X_n<1$ for all $n\ge 4$.
\end{lemma}

\begin{proof}
$X_n<1$ is $u_n<1-f_n=\frac{8n-6}{2n(2n-1)}$, i.e.\
$4\cdot 2n(2n-1)<5(n-2)(8n-6)$, i.e.\ $24n^{2}-102n+60>0$, which holds
for $n\ge 4$.
\end{proof}

\begin{lemma}\label{lem:A2}
$Y_n<1$ for all $n\ge 6$.
\end{lemma}

\begin{proof}
Since $Y_n=f_{n+1}/\bigl((1-u_n)(1-v_n)\bigr)$ and
$(1-u_n)(1-v_n)>1-u_n-v_n$, it suffices that
$u_n+v_n\le 1-f_{n+1}=\frac{8n+2}{(2n+2)(2n+1)}$. Clearing the
positive denominators, this is
$(16n-21)(4n^{2}+6n+2)\le 20(8n+2)(n-2)^{2}$, i.e.\
$96n^{3}-612n^{2}+574n+202\ge 0$, which holds for $n\ge 6$.
\end{proof}

\subsection*{Exact determinant factorization}

\begin{lemma}\label{lem:A3}
With $x=q_n$, $y=q_{n+1}$, $z=q_{n+2}$,
\[
\frac{D_n}{a_n^{3}}
=(1-y)^{2}-y^{2}(1-x)(1-z)
=(1-y)\,H(x,y)+y^{2}(1-x)(z-y),
\qquad
H(x,y):=1-y-y^{2}(1-x).
\]
In particular, by $q$-monotonicity $z\ge y$,
\[
\frac{D_n}{a_n^{3}}\ \ge\ (1-y)\,H(x,y).
\]
\end{lemma}

\begin{proof}
Expressing the nine entries of $D_n$ through $a_n$ and the ratios and
expanding gives $D_n/a_n^{3}=1-2y+xy^{2}+(1-x)y^{2}z$; both displayed
forms are elementary rearrangements. (The identity is also replayed
symbolically in the artifact and was checked on random rational
sequences.)
\end{proof}

\subsection*{Coupled monotonicity}

The function $H$ is \emph{increasing} in $x$, so substituting the
upper envelope $X_n$ for $x$ directly in $H$ would be illegitimate.
The correct substitution couples the two variables through the ratio.

\begin{lemma}\label{lem:A4}
Let $\widetilde H(x,s):=H(x,sx)=1-sx-s^{2}x^{2}(1-x)$. Then
\[
\frac{\partial\widetilde H}{\partial s}=-x\bigl(1+2sx(1-x)\bigr)<0,
\qquad
\frac{\partial\widetilde H}{\partial x}=s\bigl(-1+sx(3x-2)\bigr),
\]
and since $3x-2\le x$ for $x\le 1$, one has $sx(3x-2)\le sx^{2}<1$
whenever $sx^{2}\le S_nX_n^{2}\le Y_n<1$; hence both partial
derivatives are negative there. Consequently, if $q_n\le X_n$ and
$q_{n+1}/q_n\le S_n$ with $Y_n=S_nX_n<1$, then
\[
H(q_n,q_{n+1})=\widetilde H\bigl(q_n,\tfrac{q_{n+1}}{q_n}\bigr)
\ \ge\ \widetilde H(X_n,S_n)=H(X_n,Y_n).
\]
\end{lemma}

\begin{proof}
The two derivative identities are direct computations (the first is
replayed symbolically in the artifact). Monotone descent in $s$ then in
$x$ gives the inequality; all intermediate points satisfy the domain
condition $sx^{2}<1$ because $x\le X_n$, $s\le S_n$ and
$S_nX_n^{2}\le Y_n<1$ by Lemma~\ref{lem:A2}.
\end{proof}

\begin{lemma}\label{lem:A5}
$H(X_n,Y_n)>0$ for every $n\ge 127$.
\end{lemma}

$H(X_n,Y_n)$ is an explicit rational function.  More precisely, with
$m=n-127$,
\[
 H(X_n,Y_n)=
 \frac{P(m)}
 {(n+1)^2(2n+1)^2(5n-14)^3(20n^2-80n+69)^2},
\]
where
\begin{align*}
P(m)={}&66000m^9+71861600m^8+34728981600m^7
       +9776538943040m^6\\
 &+1766546204164425m^5+212447203445671919m^4\\
 &+17002045136190509679m^3+872992942774568354661m^2\\
 &+26091652886219863549980m
       +345765753375849332067600.
\end{align*}
Every coefficient of $P$ is strictly positive and every denominator
factor is positive for $n\ge127$, which proves the lemma directly.
The identity is independently replayed in the artifact.  Sample exact
values are $H(X_{127},Y_{127})=1.3878\ldots\times10^{-5}$ and
$H(X_{200},Y_{200})=6.566\ldots\times10^{-6}$ (the true size is
$\asymp n^{-2}$).

Combining Lemmas~\ref{lem:A2}--\ref{lem:A5}, for $n\ge127$ we have
\[
\frac{D_n}{a_n^{3}}
\ge(1-q_{n+1})H(q_n,q_{n+1})
\ge(1-Y_n)H(X_n,Y_n)>0.
\]
This is the tail closure used in \S\ref{sec:inputs}.

\begin{remark}\label{rem:sharp-envelopes}
The certificate of Lemma~\ref{lem:A5} genuinely lives at second order:
$1-Y_n$ and $Y_n^{2}(1-X_n)$ agree at order $n^{-1}$, and
$H(X_n,Y_n)\asymp n^{-2}$. Two nearby traps are worth recording.
First, the envelope \emph{constants} matter: doubling the exponent
budgets (as happens if one miscounts the step normalization of
$\kappa$) makes $H(X_n,Y_n)$ \emph{negative} for all large $n$
(already $-1.77\times10^{-4}$ at $n=127$), so no certificate can
exist. Second, the \emph{coupling} of Lemma~\ref{lem:A4} is essential:
$H$ is increasing in its first argument, and an uncoupled substitution
of upper envelopes is invalid. An earlier draft of this appendix fell
into both traps at once; we record them explicitly since either alone
silently destroys the tail closure.
\end{remark}

\begin{thebibliography}{99}

\bibitem{ASW52} M.~Aissen, I.~J.~Schoenberg, A.~M.~Whitney,
\emph{On the generating functions of totally positive sequences~I},
J. Analyse Math. \textbf{2} (1952), 93--103.

\bibitem{Edrei52} A.~Edrei,
\emph{On the generating functions of totally positive sequences~II},
J. Analyse Math. \textbf{2} (1952), 104--109.

\bibitem{CNV86} G.~Csordas, T.~S.~Norfolk, R.~S.~Varga,
\emph{The Riemann hypothesis and the Tur\'an inequalities},
Trans. Amer. Math. Soc. \textbf{296} (1986), 521--541.

\bibitem{DL11} D.~K.~Dimitrov, F.~R.~Lucas,
\emph{Higher order Tur\'an inequalities for the Riemann $\xi$-function},
Proc. Amer. Math. Soc. \textbf{139} (2011), 1013--1022.

\bibitem{GORZ19} M.~Griffin, K.~Ono, L.~Rolen, D.~Zagier,
\emph{Jensen polynomials for the Riemann zeta function and other
sequences}, Proc. Natl. Acad. Sci. USA \textbf{116} (2019),
11103--11110.

\bibitem{GORTTW} M.~Griffin, K.~Ono, L.~Rolen, J.~Thorner, Z.~Tripp,
I.~Wagner, \emph{Jensen polynomials for the Riemann xi function},
arXiv:1910.01227.

\bibitem{OSullivan} C.~O'Sullivan,
\emph{Zeros of Jensen polynomials and asymptotics for the Riemann xi
function}, arXiv:2007.13582.

\bibitem{PolyaCollected} G.~P\'olya,
\emph{\"Uber die algebraisch-funktionentheoretischen Untersuchungen von
J.~L.~W.~V.~Jensen}, Kgl. Danske Vid. Sel. Math.-Fys. Medd. \textbf{7}
(1927), 3--33.

\bibitem{Fekete13} M.~Fekete, G.~P\'olya,
\emph{\"Uber ein Problem von Laguerre},
Rend. Circ. Mat. Palermo \textbf{34} (1912), 89--120.

\bibitem{Karlin68} S.~Karlin, \emph{Total Positivity}, Vol.~I,
Stanford University Press, 1968.

\bibitem{Pinkus10} A.~Pinkus, \emph{Totally Positive Matrices},
Cambridge Tracts in Mathematics \textbf{181}, Cambridge University
Press, 2010.

\bibitem{Hutchinson23} J.~I.~Hutchinson,
\emph{On a remarkable class of entire functions},
Trans. Amer. Math. Soc. \textbf{25} (1923), 325--332.

\bibitem{RodgersTao} B.~Rodgers, T.~Tao,
\emph{The de Bruijn--Newman constant is non-negative},
Forum Math. Pi \textbf{8} (2020), e6.

\bibitem{Nuttall11} J.~Nuttall,
\emph{Determinantal approach to a proof of the Riemann hypothesis},
arXiv:1111.1128.

\bibitem{Michalowski26} W.~Micha{\l}owski,
\emph{On the P\'olya frequency order of the de Bruijn--Newman kernel:
certified failure at order five and the Toeplitz threshold phenomenon},
arXiv:2602.20313.

\bibitem{Farmer} D.~W.~Farmer,
\emph{Jensen polynomials are not a plausible route to proving the
Riemann Hypothesis}, arXiv:2008.07206.

\bibitem{PlattTrudgian} D.~Platt, T.~Trudgian,
\emph{The Riemann hypothesis is true up to $3\cdot10^{12}$},
Bull. Lond. Math. Soc. \textbf{53} (2021), 792--797.

\bibitem{Johansson17} F.~Johansson,
\emph{Arb: efficient arbitrary-precision midpoint-radius interval
arithmetic}, IEEE Trans. Computers \textbf{66} (2017), 1281--1292.

\end{thebibliography}

\end{document}
